% SPDX-License-Identifier: GPL-3.0-or-later
% Copyright (C) 2026 Alessandro Greco
\chapter{Модель угроз Tor}\label{cap:04-modello-minaccia}
\aggiornatoal{23 сентября 2026 г.}{Dingledine, Mathewson, Syverson (2004), раздел 3}

\begin{obiettivi}
К концу этой главы вы сможете:
\begin{itemize}[nosep]
  \item объяснять, что такое модель угроз и зачем она нужна каждой системе безопасности;
  \item перечислять цели и \enquote{нецели}, заявленные разработчиками Tor;
  \item описывать противника, от которого Tor должен защищать, и противника, которому он не способен противостоять;
  \item отличать то, что защищает Tor, от того, что остаётся ответственностью пользователей и приложений.
\end{itemize}
\end{obiettivi}

Ни одна система безопасности не защищает от всего. \textbf{Модель угроз} явно описывает, кого или что должна защищать система, каким противникам и угрозам она должна противостоять и какие предположения при этом делает. Именно по этому стандарту следует оценивать систему: успешная атака, укладывающаяся в модель, указывает на дефект, тогда как атака, прямо исключённая из модели, --- нет. Однако и последняя остаётся ограничением, которое пользователь обязан понимать. В этой главе рассматривается модель угроз Tor в формулировке разработчиков 2004 года~\cite{dingledine2004tor} и её последующая эволюция.

% ---------------------------------------------------------------------------
\section{Цели}\label{sec:minaccia-obiettivi}

Главная цель --- не позволить противнику \textbf{связать} две стороны одной коммуникации либо несколько коммуникаций одного пользователя. Архитектуру направляли четыре критерия, сформулированные в статье 2004 года:
\begin{description}
  \item[Развёртываемость] Система должна работать в реальном мире: не обходиться добровольцам слишком дорого, не подвергать операторов чрезмерной юридической ответственности, не требовать изменений операционной системы и не вынуждать веб-сайты что-либо устанавливать.
  \item[Удобство использования] Сложная система привлекает меньше пользователей, а меньше пользователей означает меньшую анонимность. Как объясняют авторы, удобство --- не просто комфорт, а \enquote{требование безопасности}.
  \item[Гибкость] Протокол должен быть достаточно хорошо специфицирован, чтобы служить основой для будущих исследований.
  \item[Простота] Следует использовать только хорошо изученные и проверенные методы. Каждая дополнительная функция усложняет анализ безопасности.
\end{description}

% ---------------------------------------------------------------------------
\section{Что не является целью}\label{sec:minaccia-nonobiettivi}

Не менее важно понимать, чего Tor прямо \emph{не} пытается делать:
\begin{description}
  \item[Это не одноранговая сеть (peer-to-peer)] Tor не пытается работать с тысячами краткоживущих узлов, многие из которых могут быть враждебными. Он использует набор реле, перечисленных и проверенных доверенными серверами каталогов (directory authorities) (\cref{cap:12-directory}).
  \item[Он не противостоит атакам на обоих концах соединения (end-to-end)] Tor не заявляет, что способен победить временную корреляцию или атаки пересечения, если противник наблюдает обе конечные точки.
  \item[Он не нормализует протоколы] Tor переносит данные приложений такими, каковы они есть. Если приложение раскрывает личность пользователя, Tor не может этому помешать. Изначально для этого требовался фильтрующий прокси; сегодня эту роль выполняет Tor Browser (\cref{cap:23-browser}).
  \item[Он не является стеганографическим] Tor не скрывает сам факт подключения пользователя к сети. Разработанные позднее подключаемые транспорты смягчают это ограничение (\cref{cap:16-censura}).
\end{description}

% ---------------------------------------------------------------------------
\section{Противник}\label{sec:minaccia-avversario}

Во многих теоретических исследованиях систем анонимности предполагается \textbf{глобальный пассивный противник}: тот, кто способен наблюдать весь интернет-трафик. Ни одна практическая низколатентная система не защищает от такого противника, и Tor прямо это признаёт. Вместо этого Tor предполагает противника лишь с частичной видимостью, но в некоторых отношениях с более широкими активными возможностями~\cite{dingledine2004tor}. Он может:
\begin{itemize}
  \item наблюдать \emph{часть} сетевого трафика;
  \item генерировать, изменять, удалять или задерживать трафик;
  \item эксплуатировать собственные реле;
  \item компрометировать некоторые реле, управляемые другими.
\end{itemize}
С этими возможностями противник может пытаться проводить множество атак: коррелировать входящий и исходящий трафик, наблюдая края сети; компрометировать реле или ключи; повторно передавать трафик; выводить честные реле из строя, чтобы направлять пользователей к вредоносным (\cref{cap:15-dos}); маркировать трафик так, чтобы распознавать его в другом месте (\cref{sec:tor1-debolezze}); через систему каталогов показывать разным пользователям несовместимые представления сети (\cref{cap:12-directory}); либо сделать сеть настолько ненадёжной, чтобы пользователи отказались от неё в пользу менее защищённых систем.

\subsection{Подтверждение и анализ}

Разработчики Tor различают два типа атак на основе трафика. При атаке \textbf{подтверждения} противник уже подозревает наличие связи (\enquote{Алиса общается с Бобом}) и наблюдает обе конечные точки, чтобы проверить гипотезу. Tor, как и любая низколатентная система, не способен полностью защищать от таких атак. При атаке \textbf{анализа} противник изучает трафик, чтобы выяснить, \emph{где} наблюдать, то есть какие точки сети следует атаковать. В первую очередь Tor рассчитан на сопротивление именно этому классу. Весь процесс выбора маршрута, включая guard-реле и весовые коэффициенты пропускной способности (\cref{cap:13-percorso}), призван снизить вероятность того, что противник окажется одновременно на обоих концах цепи.

\begin{attenzione}[Главный принцип]
Tor защищает от того, кто наблюдает \emph{часть} сети, но не от того, кому удаётся наблюдать \emph{оба конца} коммуникации. Противник, контролирующий или наблюдающий одновременно соединение пользователя с guard-реле и соединение exit-реле с веб-сайтом, может коррелировать временные характеристики и объём трафика. Для организации, имеющей доступ к сетям целой страны, и пользователь, и веб-сайт могут находиться внутри этой страны. Это практический, а не только теоретический сценарий.
\end{attenzione}

% ---------------------------------------------------------------------------
\section{От чего Tor защищает, а от чего нет}\label{sec:minaccia-sintesi}

\begin{figure}[htbp]
\centering
% SPDX-License-Identifier: GPL-3.0-or-later
% Copyright (C) 2026 Alessandro Greco
%% Русское альтернативное описание EPUB.
\altfig{Цепь Tor от клиента через сторожевой, промежуточный и выходной ретрансляторы к сайту назначения. Таблица под ретрансляторами и сайтом показывает видимость на сетевом уровне без утечек приложения: сторожевой узел видит исходный IP-адрес клиента, выходной знает назначение, а промежуточный не знает ни один из двух концов. Сайт видит IP выходного узла вместо IP клиента; наблюдатель на входе и выходе может сопоставлять потоки.}
% Что знает каждый наблюдатель вдоль цепи.
\begin{tikzpicture}[font=\sffamily\scriptsize,
  sa/.style={fill=avviso!15,draw=avviso!60,rounded corners=1pt,minimum height=0.5cm,minimum width=1.5cm},
  no/.style={fill=torverde!12,draw=torverde!50!black,rounded corners=1pt,minimum height=0.5cm,minimum width=1.5cm}]
  \node[client,minimum width=1.5cm,minimum height=0.7cm] (c) at (0,0) {клиент};
  \node[relay,minimum size=0.9cm] (g) at (2.6,0) {сторожевой};
  \node[relay,minimum size=0.9cm] (m) at (5.2,0) {промежуточный};
  \node[relay,minimum size=0.9cm,fill=torverde!30] (e) at (7.8,0) {выходной};
  \node[server,minimum width=1.5cm,minimum height=0.7cm] (s) at (10.4,0) {сайт};
  \draw[cellaf] (c) -- (g); \draw[cellaf] (g) -- (m); \draw[cellaf] (m) -- (e);
  \draw[->,thick,torgrigio] (e) -- (s);
  \node[font=\sffamily\tiny,anchor=east,align=right] at (-1.0,-1.1) {знает\\клиента};
  \node[font=\sffamily\tiny,anchor=east,align=right] at (-1.0,-1.9) {знает\\назначение};
  \node[sa] at (2.6,-1.1) {да};
  \node[no] at (5.2,-1.1) {нет}; \node[no] at (7.8,-1.1) {нет}; \node[no] at (10.4,-1.1) {нет};
  \node[no] at (2.6,-1.9) {нет}; \node[no] at (5.2,-1.9) {нет};
  \node[sa] at (7.8,-1.9) {да}; \node[sa] at (10.4,-1.9) {да};
  \node[align=center,font=\sffamily\tiny,text=avviso] at (5.2,-2.8) {ни один ретранслятор Tor не знает оба конца цепи:\\наблюдатель на входе и выходе может сопоставить потоки (гл.~26)};
\end{tikzpicture}

\caption{Что знает каждая точка маршрута. Защита Tor основана на разделении: guard-реле знает клиента, но не пункт назначения; exit-реле знает пункт назначения, но не клиента.}
\label{fig:sapere}
\end{figure}

\begin{center}\small
\begin{tabularx}{\textwidth}{@{}>{\raggedright\arraybackslash}X>{\raggedright\arraybackslash}X@{}}
\toprule
Tor защищает от & Tor не защищает от \\
\midrule
веб-сайта, пытающегося узнать IP-адрес посетителя & информации, которую сам пользователь раскрывает, например входя в аккаунт или сообщая личные сведения \\
наблюдателя в локальной сети или интернет-провайдера, видящего только соединение с Tor & наблюдателя, следящего за обоими концами коммуникации \\
отдельного реле, знающего не более одной конечной точки & противника, контролирующего одновременно guard- и exit-реле цепи \\
exit-реле, пытающегося определить отправителя запроса & exit-реле, читающего или изменяющего незашифрованный трафик без HTTPS \\
цензуры по IP-адресам при использовании мостов & уязвимого, скомпрометированного или неправильно настроенного приложения либо браузера \\
\bottomrule
\end{tabularx}
\end{center}

Со временем модель расширилась, чтобы учесть явления, которые статья 2004 года не могла предвидеть: цензуру в масштабах целых стран (\cref{cap:16-censura}), защиту от \emph{фингерпринтинга} браузера (\cref{cap:23-browser}), защиту от DoS (\cref{cap:15-dos}) и будущую угрозу квантовых компьютеров (\cref{cap:31-futuro}). Один принцип остаётся неизменным: понимание модели угроз --- первый шаг к безопасному использованию Tor.

% ---------------------------------------------------------------------------
\section*{Итоги}
\begin{itemize}
  \item Модель угроз определяет противников и атаки, от которых система предназначена защищать.
  \item Tor стремится не позволить связать пользователя с пунктом назначения, отдавая приоритет развёртываемости, удобству, гибкости и простоте.
  \item Tor не является peer-to-peer системой, не побеждает end-to-end атаки, не нормализует прикладные протоколы и без подключаемых транспортов не скрывает факт использования Tor.
  \item Противник может наблюдать часть сети, эксплуатировать и компрометировать реле и предпринимать активные действия. Tor рассчитан прежде всего на сопротивление атакам анализа и слабее защищает от атак подтверждения.
\end{itemize}

% ---------------------------------------------------------------------------
\section*{Упражнения}

\begin{esercizio}[1]\label{es:04-casi}
Для каждого сценария укажите, защищает ли Tor пользователя, и объясните почему: (a) оператор Wi-Fi в кафе хочет узнать, какие сайты посещает клиент; (b) веб-сайт хочет узнать IP-адреса посетителей; (c) пользователь через Tor входит в личный аккаунт социальной сети; (d) ведомство, контролирующее сети страны, подозревает, что пользователь общается с газетой в той же стране.
\end{esercizio}

\begin{esercizio}[2]\label{es:04-usabilita}
Приведите пример, показывающий, почему удобство использования в системе анонимности является требованием безопасности, а не просто вопросом комфорта.
\end{esercizio}

\begin{esercizio}[2]\label{es:04-modello}
Составьте модель угроз для журналиста, который использует Tor Browser для общения с источником в другой стране. Определите возможных противников, то, что они способны наблюдать, и аспекты, остающиеся за пределами защиты Tor.
\end{esercizio}

% ---------------------------------------------------------------------------
\section*{Дополнительная литература}
\begin{itemize}
  \item R.~Dingledine, N.~Mathewson, P.~Syverson, \titolo{Tor: The Second-Generation Onion Router}~\cite{dingledine2004tor}, разделы 3 и 7.
  \item R.~Dingledine, N.~Mathewson, \titolo{Anonymity Loves Company: Usability and the Network Effect}~\cite{dingledine2006company}.
\end{itemize}
